%Kína TST 1994
%combi 7
Við segjum að $n \times n$ fylki af tölum sé \emph{$n$-kóði} ef tölurnar í öllum röðum og öllum dálkum mynda jafnmunarunur. Mengi reita í fylkinu er kallað \emph{lykill}, ef með því að sýna aðeins tölurnar í þeim reitum er hægt að ótvírætt ákvarða allar aðrar tölur.
\begin{enumerate}[label=\alph*.]
\item 
Gefið að $n\ge 4$, finnið minnstu jákvæðu heiltöluna $s$, svo hvaða $s$ reitir sem er í $n$-kóða mynda lykil.
\item
Gefið að $n\ge 4 $ finnið minnstu jákvæðu töluna $t$, svo hvaða $t$ reitir sem er á skálínu $n$-kóða mynda lykil.
\end{enumerate}